\documentclass[UTF8,a4paper,11pt]{ctexart}
\usepackage{amsmath,amssymb}
\usepackage[margin=2.5cm]{geometry}
\usepackage[hidelinks]{hyperref}
\title{2023 年北京卷第 20 题：独立数学证明}
\date{2026 年 10 月 25 日}
\begin{document}
\maketitle
\noindent\textbf{固定题源：}\url{https://github.com/DxAThing/Gaokao-Math-Problems-Compilation/blob/f1ae169b58a75b3db6e35c670b3f0d5c78718188/content/2023/beijing.tex#L627}

设 $f(x)=x-x^3e^{ax+b}$ 在 $x=1$ 处的切线为 $y=-x+1$。
\paragraph{第一问}
由 $f(1)=0,\ f'(1)=-1$ 和
$f'(x)=1-(3x^2+ax^3)e^{ax+b}$，
得 $e^{a+b}=1$、$(3+a)e^{a+b}=2$，故
\[\boxed{a=-1,\qquad b=1}.\]
\paragraph{第二问}
令 $g=f'$，则
\[
g(x)=1-x^2(3-x)e^{1-x},\qquad
g'(x)=-x(x^2-6x+6)e^{1-x}.
\]
记 $\alpha=3-\sqrt3$，$\beta=3+\sqrt3$。
因为指数因子恒正，$g'$ 在
$(-\infty,0)$、$(0,\alpha)$、$(\alpha,\beta)$、
$(\beta,+\infty)$ 上的符号依次为正、负、正、负。
因此 $g$ 在第一、三段严格递增，在第二、四段严格递减。

\paragraph{第三问}
有 $g(-1)=1-4e^2<0$，$g(0)=1>0$，
$g(1)=-1<0$，$g(3)=1>0$。
因 $1<\alpha$，又 $g$ 在 $(0,\alpha)$ 严格递减，
故 $g(\alpha)<g(1)<0$。结合连续性、严格单调性，
$g$ 在 $(-1,0)$、$(0,1)$、$(\alpha,3)$
上各有且仅有一个零点，其符号分别由负变正、由正变负、
由负变正；各单调区间内不可能出现额外零点。
且当 $x\ge3$ 时，$x^2(3-x)\le0$，从而 $g(x)\ge1>0$。
综上全部实数区间已覆盖，函数 $f$ 恰有
两个极小值点、一个极大值点：
\[
\boxed{\text{极值点共 }3\text{ 个}}.
\]

\noindent\textbf{审核状态：}数学推导已独立核验；
未进行 XeLaTeX 编译或 PDF 视觉质检。
\end{document}
